\begin{acknowledgements}
	A mi familia, no solo a mis padres sino a cada uno de los que estuvieron cerca, por sostener con paciencia cada etapa de la carrera y celebrar los avances como propios.
	
	A mis amigos, por recordarme en los momentos justos que hay vida más allá de la facultad, y por comprender mis ausencias durante los tramos más exigentes del trabajo.
	
	A los buenos profesores que tuvimos a lo largo de la formación, los que enseñaron con rigor y vocación y marcaron la diferencia entre asistir a clase y aprender de verdad. Su enseñanza vive en cada decisión técnica del proyecto, aunque sus nombres no aparezcan en la bibliografía.
	
	A mis tutores, por adoptarnos cuando este proyecto todavía era un esquema impreciso y acompañarlo paso a paso hasta convertirlo en lo que el lector tiene ahora en las manos.
	
	Al Instituto Cubano de Investigación Cultural Juan Marinello, por confiar en que un trabajo de fin de carrera podía aportarle valor concreto a la digitalización de su acervo y por abrirnos las puertas para llevarlo a cabo.
	
	A quienes se sientan identificados leyendo estas lineas, muchas gracias a todos, ustedes son los pilares sobre los cuales este proyecto se hizo realidad.
	
\end{acknowledgements}